\chapter{แนวคิดและทฤษฎีพื้นฐานของการเรียนรู้แบบมีผู้สอน}

\section{หลักการและแนวคิดของการเรียนรู้แบบมีผู้สอน}
การเรียนรู้แบบมีผู้สอน (Supervised Learning) เป็นสาขาหนึ่งของการเรียนรู้ของเครื่อง (Machine Learning) ที่มีบทบาทสำคัญในการสร้างระบบที่สามารถทำนายผลลัพธ์จากข้อมูลที่ได้รับการระบุคำตอบไว้ล่วงหน้าแล้ว หลักการพื้นฐานของการเรียนรู้แบบมีผู้สอนคือการที่ระบบคอมพิวเตอร์ได้รับตัวอย่างข้อมูลที่ประกอบด้วยทั้งข้อมูลนำเข้าและคำตอบที่ถูกต้อง (Labeled Data) เพื่อเรียนรู้รูปแบบความสัมพันธ์ระหว่างข้อมูลนำเข้าและผลลัพธ์ และนำความรู้ที่ได้ไปประยุกต์ใช้กับข้อมูลใหม่ที่ไม่เคยพบเห็นมาก่อน กระบวนการนี้เปรียบเสมือนการที่ครูสอนนักเรียนโดยการยกตัวอย่างพร้อมเฉลยคำตอบ แล้วให้นักเรียนจดจำและเข้าใจรูปแบบเพื่อตอบคำถามในข้อสอบที่ยังไม่เคยเห็นได้อย่างถูกต้อง

ในทางคณิตศาสตร์ การเรียนรู้แบบมีผู้สอนสามารถอธิบายได้ดังนี้ สมมติให้ $\mathcal{D} = \{(\mathbf{x}_i, y_i)\}_{i=1}^{n}$ แทนชุดข้อมูลฝึกสอน (Training Dataset) โดยที่ $\mathbf{x}_i \in \mathbb{R}^{d}$ แทนเวกเตอร์ของคุณลักษณะ (Feature Vector) ของตัวอย่างที่ $i$ และ $y_i$ แทนค่าคำตอบหรือฉลาก (Label) ที่สอดคล้องกับ $\mathbf{x}_i$ โดยที่ $y_i$ อาจเป็นค่าต่อเนื่อง (Continuous) ในกรณีของการถดถอย (Regression) หรือเป็นค่าไม่ต่อเนื่อง (Discrete) ในกรณีของการจำแนกประเภท (Classification) วัตถุประสงค์ของการเรียนรู้แบบมีผู้สอนคือการค้นหาฟังก์ชัน $f: \mathbb{R}^{d} \rightarrow \mathcal{Y}$ ที่สามารถอธิบายความสัมพันธ์ระหว่าง $\mathbf{x}$ และ $y$ ได้อย่างแม่นยำ โดยฟังก์ชัน $f$ นี้จะอยู่ในรูปของตัวแบบ (Model) หรือตัวประมาณ (Estimator) ที่มีพารามิเตอร์ $\boldsymbol{\theta}$ ซึ่งต้องได้รับการเรียนรู้จากข้อมูลฝึกสอน

ความแตกต่างที่สำคัญระหว่างการเรียนรู้แบบมีผู้สอนกับการเรียนรู้แบบไม่มีผู้สอน (Unsupervised Learning) อยู่ที่การมีฉลาก (Label) กำกับข้อมูล ในการเรียนรู้แบบไม่มีผู้สอน ระบบจะได้รับข้อมูลที่ไม่มีคำตอบกำกับและต้องค้นหาโครงสร้างหรือรูปแบบที่ซ่อนอยู่ด้วยตัวเอง ตัวอย่างเช่น การแบ่งกลุ่ม (Clustering) หรือการลดมิติ (Dimensionality Reduction) ในขณะที่การเรียนรู้แบบมีผู้สอนมีคำตอบที่ถูกต้องเป็นเป้าหมายในการเรียนรู้ ทำให้สามารถวัดประสิทธิภาพของตัวแบบได้โดยตรงจากความผิดพลาดระหว่างค่าทำนายและค่าจริง

การเรียนรู้แบบมีผู้สอนสามารถจำแนกประเภทหลักออกได้เป็นสองกลุ่มใหญ่ตามลักษณะของตัวแปรตาม กลุ่มแรกคือการถดถอย (Regression) ซึ่งตัวแปรตามเป็นข้อมูลต่อเนื่อง เช่น การทำนายราคาบ้าน การพยากรณ์อุณหภูมิ หรือการประมาณค่ายอดขาย กลุ่มที่สองคือการจำแนกประเภท (Classification) ซึ่งตัวแปรตามเป็นข้อมูลไม่ต่อเนื่องหรือเชิงคุณภาพ เช่น การแยกอีเมลขยะ การวินิจฉัยโรค หรือการจดจำใบหน้า ทั้งสองกลุ่มนี้มีหลักการพื้นฐานเหมือนกันคือการเรียนรู้การจับคู่ระหว่างข้อมูลนำเข้าและผลลัพธ์ แต่แตกต่างกันในรายละเอียดของฟังก์ชันวัตถุประสงค์และวิธีการประเมินประสิทธิภาพ

\begin{myBox}[]{องค์ประกอบหลักของการเรียนรู้แบบมีผู้สอน}{}
การเรียนรู้แบบมีผู้สอนประกอบด้วยองค์ประกอบสำคัญ 4 ประการที่ต้องพิจารณาในการสร้างตัวแบบที่มีประสิทธิภาพ ดังนี้

\begin{enumerate}
\item \textbf{ข้อมูลฝึกสอน (Training Data)} คือชุดข้อมูลที่ประกอบด้วยคู่ลำดับ $(\mathbf{x}_i, y_i)$ จำนวน $n$ ตัวอย่าง โดยที่ $\mathbf{x}_i$ เป็นเวกเตอร์ของคุณลักษณะและ $y_i$ คือฉลากที่ถูกต้อง

\item \textbf{ตัวแบบ (Model)} คือฟังก์ชัน $f(\mathbf{x}; \boldsymbol{\theta})$ ที่มีพารามิเตอร์ $\boldsymbol{\theta}$ ซึ่งต้องเรียนรู้จากข้อมูล เช่น น้ำหนักในโครงข่ายประสาทเทียม หรือสัมประสิทธิ์ในสมการถดถอย

\item \textbf{ฟังก์ชันวัตถุประสงค์ (Objective Function)} คือฟังก์ชันสูญเสีย (Loss Function) $\mathcal{L}(y, \hat{y})$ ที่วัดความผิดพลาดระหว่างค่าจริง $y$ และค่าทำนาย $\hat{y} = f(\mathbf{x}; \boldsymbol{\theta})$

\item \textbf{ขั้นตอนการหาค่าเหมาะสมที่สุด (Optimization)} คือกระบวนการปรับพารามิเตอร์ $\boldsymbol{\theta}$ ให้ลดค่าฟังก์ชันสูญเสียให้น้อยที่สุดบนข้อมูลฝึกสอน
\end{enumerate}
\end{myBox}

ในทางปฏิบัติ การเรียนรู้แบบมีผู้สอนถูกนำไปประยุกต์ใช้อย่างกว้างขวางในหลากหลายสาขา ตัวอย่างที่พบได้บ่อยในชีวิตประจำวัน เช่น ระบบคัดกรองอีเมลขยะ (Spam Filtering) ที่เรียนรู้จากอีเมลที่ผู้ใช้ระบุว่าเป็นขยะหรือไม่เป็นขยะ ระบบแนะนำสินค้า (Product Recommendation) ที่เรียนรู้จากประวัติการซื้อและการเข้าชมของลูกค้า ระบบตรวจจับการฉ้อโกง (Fraud Detection) ที่เรียนรู้จากธุรกรรมที่เป็นการฉ้อโกงและไม่เป็นการฉ้อโกง ระบบวินิจฉัยโรค (Medical Diagnosis) ที่เรียนรู้จากผลตรวจและการวินิจฉัยของแพทย์ และระบบพยากรณ์ราคา (Price Forecasting) ที่เรียนรู้จากข้อมูลราคาในอดีตและปัจจัยที่เกี่ยวข้อง ความสำเร็จของการประยุกต์ใช้เหล่านี้ขึ้นอยู่กับคุณภาพของข้อมูลฝึกสอน การออกแบบคุณลักษณะ (Feature Engineering) ที่เหมาะสม และการเลือกวิธีการเรียนรู้ที่เหมาะสมกับลักษณะของปัญหา

\begin{exmp}
\label{exmp:supervised_learning_iris}
พิจารณาตัวอย่างการสร้างระบบจำแนกประเภทดอกไม้ไอริส (Iris Dataset) ซึ่งเป็นชุดข้อมูลมาตรฐานในการเรียนรู้ของเครื่อง ชุดข้อมูลประกอบด้วยดอกไม้ 150 ตัวอย่าง แต่ละตัวอย่างมี 4 คุณลักษณะ ได้แก่ ความยาวกลีบเลี้ยง (Sepal Length) ความกว้างกลีบเลี้ยง (Sepal Width) ความยาวกลีบดอก (Petal Length) และความกว้างกลีบดอก (Petal Width) และมี 3 ชนิด ได้แก่ Setosa, Versicolor, และ Virginica

ขั้นตอนการสร้างระบบจำแนกประเภทด้วยการเรียนรู้แบบมีผู้สอนมีดังนี้

\textbf{ขั้นตอนที่ 1: เตรียมข้อมูลฝึกสอน} แบ่งข้อมูล 150 ตัวอย่างออกเป็นชุดฝึกสอน 120 ตัวอย่าง (80\%) และชุดทดสอบ 30 ตัวอย่าง (20\%) โดยใช้การสุ่มแบ่งแบบ Stratified เพื่อรักษาสัดส่วนของแต่ละชนิดดอกไม้ให้เท่าเดิม

\textbf{ขั้นตอนที่ 2: เลือกตัวแบบ} เลือกใช้ Naive Bayes Classifier เนื่องจากข้อมูลมีลักษณะเป็นข้อมูลต่อเนื่อง และสมมติว่าคุณลักษณะแต่ละตัวมีการแจกแจงแบบปกติ (Gaussian Distribution)

\textbf{ขั้นตอนที่ 3: ฝึกสอนตัวแบบ} คำนวณค่าเฉลี่ย ($\mu$) และส่วนเบี่ยงเบนมาตรฐาน ($\sigma$) ของแต่ละคุณลักษณะสำหรับแต่ละชนิดดอกไม้จากชุดข้อมูลฝึกสอน

\textbf{ขั้นตอนที่ 4: ประเมินประสิทธิภาพ} นำตัวแบบไปทดสอบกับชุดข้อมูลทดสอบ 30 ตัวอย่างที่ไม่เคยเห็นมาก่อน สมมติว่าได้ผลการทำนายดังนี้: ทำนายถูกต้อง 28 ตัวอย่าง ทำนายผิด 2 ตัวอย่าง ให้ค่า Accuracy $= \frac{28}{30} \times 100 \approx 93.3\%$

สรุปได้ว่าระบบจำแนกประเภทดอกไม้ไอริสที่สร้างขึ้นสามารถจำแนกชนิดดอกไม้ได้ถูกต้องประมาณ 93.3\% เมื่อพบข้อมูลใหม่ที่ไม่เคยเห็นมาก่อน
\end{exmp}

\section{ประเภทของปัญหาการจำแนกประเภทและการพยากรณ์}
การจำแนกประเภท (Classification) และการพยากรณ์ (Prediction) เป็นหัวใจสำคัญของการเรียนรู้แบบมีผู้สอน แม้ว่าทั้งสองจะมีเป้าหมายในการทำนายผลลัพธ์จากข้อมูลนำเข้า แต่ความแตกต่างพื้นฐานอยู่ที่ลักษณะของผลลัพธ์ที่ต้องการทำนาย การพยากรณ์มักหมายถึงการทำนายค่าต่อเนื่อง (Continuous Value) ซึ่งเป็นข้อมูลเชิงปริมาณที่สามารถรับค่าใดก็ได้ในช่วงจำนวนจริงหนึ่ง ในขณะที่การจำแนกประเภทมุ่งเน้นการจัดข้อมูลเข้าเป็นหมวดหมู่หรือกลุ่มที่กำหนดไว้ล่วงหน้า ซึ่งเป็นข้อมูลเชิงคุณภาพที่มีจำนวนจำกัดของค่า

ในการพยากรณ์ ผลลัพธ์ $y$ ที่ต้องการทำนายเป็นตัวแปรต่อเนื่องที่มีค่าในช่วง $(-\infty, \infty)$ หรือช่วงจำนวนจริงบวกในกรณีของตัวแปรที่ไม่เป็นลบ ตัวอย่างเช่น การพยากรณ์ราคาบ้านซึ่งอาจมีค่าเป็น 2.5 ล้านบาท หรือ 4.3 ล้านบาท การพยากรณ์อุณหภูมิซึ่งอาจเป็น 25.6 องศาเซลเซียส หรือการพยากรณ์ยอดขายซึ่งอาจเป็น 1.2 ล้านบาท ลักษณะต่อเนื่องนี้ทำให้สามารถวัดความผิดพลาดในเชิงปริมาณได้ เช่น ความแตกต่างระหว่างค่าจริงและค่าทำนายมีหน่วยเดียวกันกับตัวแปรตาม

ในทางคณิตศาสตร์ ปัญหาการพยากรณ์สามารถเขียนได้ว่าต้องการหาฟังก์ชัน $f: \mathbb{R}^{d} \rightarrow \mathbb{R}$ ที่ทำการทำแมปเวกเตอร์คุณลักษณะ $\mathbf{x} \in \mathbb{R}^{d}$ ไปยังค่าต่อเนื่อง $y \in \mathbb{R}$ โดยมีเกณฑ์ความสำเร็จคือค่าทำนาย $\hat{y} = f(\mathbf{x})$ ควรใกล้เคียงกับค่าจริง $y$ มากที่สุดเท่าที่จะเป็นไปได้ การวัดความใกล้เคียงนี้มักใช้ฟังก์ชันสูญเสีย เช่น Mean Squared Error (MSE) ซึ่งมีสูตร $\mathcal{L}(y, \hat{y}) = (y - \hat{y})^{2}$ หรือ Mean Absolute Error (MAE) ซึ่งมีสูตร $\mathcal{L}(y, \hat{y}) = |y - \hat{y}|$ การเลือกฟังก์ชันสูญเสียขึ้นอยู่กับลักษณะของปัญหาและคุณสมบัติที่ต้องการของตัวประมาณ

สำหรับการจำแนกประเภท ผลลัพธ์ $y$ ที่ต้องการทำนายเป็นตัวแปรไม่ต่อเนื่องที่แทนหมวดหมู่หรือกลุ่ม กลุ่มเหล่านี้อาจเป็นค่าข้อความ (เช่น "แม่กับลูกสุนัข" หรือ "แม่กับลูกแมว") หรือค่าตัวเลข (เช่น 0 หรือ 1 สำหรับการจำแนกแบบทวิภาค) หรือค่าหลายกลุ่ม (เช่น 1, 2, 3 สำหรับการจำแนกหลายกลุ่ม) ในกรณีทวิภาค (Binary Classification) ค่าฉลากมักเขียนแทนด้วย $y \in \{0, 1\}$ โดยที่ 0 แทนการไม่เกิดเหตุการณ์ที่สนใจ (เช่น ไม่เป็นโรค ไม่ฉ้อโกง ไม่ผิดนัด) และ 1 แทนการเกิดเหตุการณ์ที่สนใจ (เช่น เป็นโรค ฉ้อโกง ผิดนัด) ในกรณีหลายกลุ่ม (Multi-class Classification) ค่าฉลากเขียนแทนด้วย $y \in \{1, 2, \ldots, K\}$ โดยที่ $K$ คือจำนวนกลุ่ม

การจำแนกประเภทสามารถแบ่งย่อยได้เป็นหลายประเภทตามลักษณะของปัญหา ประเภทแรกคือการจำแนกแบบทวิภาค (Binary Classification) ซึ่งมีเพียงสองกลุ่มเท่านั้น ตัวอย่างเช่น การตัดสินใจอนุมัติสินเชื่อ (อนุมัติหรือไม่อนุมัติ) การวินิจฉัยโรค (เป็นโรคหรือไม่เป็น) หรือการตรวจจับสแปม (เป็นสแปมหรือไม่เป็น) ประเภทที่สองคือการจำแนกหลายกลุ่ม (Multi-class Classification) ซึ่งมีมากกว่าสองกลุ่ม ตัวอย่างเช่น การจดจำตัวอักษร (A, B, C, \ldots, Z) การจำแนกประเภทผลไม้ (ส้ม แอปเปิ้ล กล้วย มะม่วง) หรือการจำแนกระดับความเสี่ยง (ต่ำ กลาง สูง สูงมาก) ประเภทที่สามคือการจำแนกหลายป้ายกำกับ (Multi-label Classification) ซึ่งข้อมูลแต่ละตัวอย่างอาจเป็นสมาชิกของหลายกลุ่มพร้อมกัน ตัวอย่างเช่น การติดแท็กบทความ (บทความหนึ่งอาจมีแท็ก "การเมือง" "เศรษฐกิจ" และ "สังคม" พร้อมกัน)

\begin{myBox}[]{ตัวอย่างการประยุกต์ใช้การจำแนกประเภทและการพยากรณ์}{}
เพื่อให้เข้าใจความแตกต่างระหว่างการจำแนกประเภทและการพยากรณ์อย่างชัดเจน พิจารณาตัวอย่างต่อไปนี้

สมมติมีข้อมูลลูกค้าธนาคารและต้องการทำนายผลลัพธ์ที่แตกต่างกัน หากต้องการทำนายว่าลูกค้าจะผิดนัดชำระหนี้หรือไม่ นี่คือปัญหาการจำแนกประเภทแบบทวิภาค โดยผลลัพธ์เป็น $y \in \{0, 1\}$ คือ "ไม่ผิดนัด" หรือ "ผิดนัด" หากต้องการทำนายจำนวนเดือนที่ลูกค้าจะผิดนัดชำระหนี้ (เช่น 0, 1, 2, 3 เดือน) นี่ก็ยังคงเป็นการจำแนกประเภทแบบหลายกลุ่ม แต่หากต้องการทำนายมูลค่าความเสียหายที่อาจเกิดขึ้นเป็นจำนวนเงินที่เป็นค่าต่อเนื่อง นี่คือปัญหาการพยากรณ์

ความแตกต่างสำคัญอีกประการหนึ่งคือการตีความผลลัพธ์ ในการจำแนกประเภท ผลลัพธ์ที่ได้คือการจัดกลุ่มหรือความน่าจะเป็นของการเป็นสมาชิกแต่ละกลุ่ม ในขณะที่การพยากรณ์ให้ผลลัพธ์เป็นค่าตัวเลขที่มีหน่วยเดียวกันกับตัวแปรตาม การเลือกประเภทปัญหาที่ถูกต้องขึ้นอยู่กับลักษณะของคำถามทางธุรกิจและวัตถุประสงค์ของการวิเคราะห์
\end{myBox}

ในการประยุกต์ใช้จริง การตัดสินใจว่าปัญหาที่สนใจเป็นการจำแนกประเภทหรือการพยากรณ์ขึ้นอยู่กับลักษณะของคำถามทางธุรกิจและวัตถุประสงค์ในการใช้ผลลัพธ์ บางครั้งปัญหาเดียวกันอาจมองได้ทั้งสองมุมมอง ตัวอย่างเช่น การทำนายราคาหุ้นอาจมองได้ทั้งเป็นการพยากรณ์ (ทำนายราคาที่แน่นอน) หรือเป็นการจำแนกประเภท (ทำนายว่าราคาจะขึ้นหรือลง) การเลือกมุมมองที่เหมาะสมขึ้นอยู่กับความต้องการในการใช้งานและความถูกต้องของข้อมูล

สำหรับปัญหาการจำแนกประเภทหลายกลุ่ม มีความซับซ้อนเพิ่มขึ้นเนื่องจากต้องพิจารณาความสัมพันธ์ระหว่างกลุ่มต่างๆ ตัวแบบที่ใช้อาจออกแบบให้ทำนายกลุ่มโดยตรง (Direct Multi-class) หรืออาจใช้การรวมตัวแบบแบบทวิภาคหลายตัวเข้าด้วยกัน เช่น แนวทาง One-vs-Rest (OvR) ที่สร้างตัวแบบทวิภาค $K$ ตัวสำหรับ $K$ กลุ่ม โดยแต่ละตัวแบบจะแยกกลุ่มหนึ่งออกจากกลุ่มที่เหลือ หรือแนวทาง One-vs-One (OvO) ที่สร้างตัวแบบทวิภาค $\frac{K(K-1)}{2}$ ตัวสำหรับทุกคู่ของกลุ่ม การเลือกแนวทางขึ้นอยู่กับลักษณะของตัวแบบพื้นฐานและประสิทธิภาพในการคำนวณ

\section{ข้อมูลฝึกสอนและข้อมูลทดสอบ}
ข้อมูลฝึกสอน (Training Data) และข้อมูลทดสอบ (Test Data) เป็นองค์ประกอบพื้นฐานที่แบ่งแยกกระบวนการเรียนรู้ออกจากการประเมินประสิทธิภาพในการเรียนรู้แบบมีผู้สอน ข้อมูลฝึกสอนคือชุดข้อมูลที่ใช้ในการเรียนรู้พารามิเตอร์ของตัวแบบ โดยตัวแบบจะปรับพารามิเตอร์ของตัวเองเพื่อลดค่าฟังก์ชันสูญเสียบนข้อมูลชุดนี้ให้น้อยที่สุด ในขณะที่ข้อมูลทดสอบคือชุดข้อมูลที่แยกออกไว้และไม่เคยถูกใช้ในกระบวนการฝึกสอน ใช้สำหรับประเมินประสิทธิภาพที่แท้จริงของตัวแบบเมื่อพบข้อมูลใหม่ที่ไม่เคยเห็นมาก่อน

ความสำคัญของการแบ่งข้อมูลออกเป็นชุดฝึกสอนและชุดทดสอบมีรากฐานจากเป้าหมายหลักของการเรียนรู้แบบมีผู้สอน นั่นคือการสร้างตัวแบบที่สามารถทำนายผลลัพธ์ได้อย่างแม่นยำสำหรับข้อมูลใหม่ที่ยังไม่เคยพบเห็น ไม่ใช่เพียงการจดจำข้อมูลฝึกสอน หากตัวแบบมีความซับซ้อนมากเกินไปจนสามารถทำนายข้อมูลฝึกสอนได้อย่างสมบูรณ์แบบ (ค่าความผิดพลาดบนชุดฝึกสอนเป็นศูนย์) แต่ไม่สามารถทำนายข้อมูลใหม่ได้ดี ปรากฏการณ์นี้เรียกว่า Overfitting ซึ่งเป็นปัญหาพื้นฐานที่ต้องหลีกเลี่ยงในการสร้างตัวแบบการเรียนรู้ของเครื่อง

สมมติให้ชุดข้อมูลทั้งหมด $\mathcal{D}$ ถูกแบ่งออกเป็นสองส่วนที่ไม่ทับซ้อนกัน คือ $\mathcal{D}_{\text{train}}$ และ $\mathcal{D}_{\text{test}}$ โดยที่ $\mathcal{D}_{\text{train}} \cap \mathcal{D}_{\text{test}} = \emptyset$ และ $\mathcal{D}_{\text{train}} \cup \mathcal{D}_{\text{test}} = \mathcal{D}$ หรืออาจเขียนในรูป $\mathcal{D} = \mathcal{D}_{\text{train}} \cup \mathcal{D}_{\text{test}}$ เมื่อทั้งสองส่วนไม่มีสมาชิกร่วมกัน การแบ่งสัดส่วนที่นิยมใช้คือ 70:30 หรือ 80:20 ของข้อมูลทั้งหมดสำหรับฝึกสอนและทดสอบตามลำดับ อย่างไรก็ตาม สัดส่วนที่เหมาะสมขึ้นอยู่กับขนาดของข้อมูลทั้งหมด หากมีข้อมูลจำนวนมาก อาจใช้สัดส่วน 90:10 หรือแม้แต่ 95:5 ก็ได้ เพราะชุดทดสอบที่มีข้อมูลจำนวนหลายพันหรือหลายหมื่นตัวอย่างก็เพียงพอสำหรับการประเมินประสิทธิภาพอย่างน่าเชื่อถือ

การแบ่งข้อมูลที่ดีควรคำนึงถึงหลายประเด็นสำคัญ ประการแรกคือการสุ่มแบ่ง (Random Splitting) ควรสุ่มตัวอย่างไปยังชุดฝึกสอนและชุดทดสอบอย่างสุ่ม ไม่ใช่การแบ่งตามลำดับของข้อมูล เพื่อป้องกันการรั่วไหลของข้อมูล (Data Leakage) ที่อาจเกิดขึ้นจากรูปแบบเฉพาะที่อยู่ในลำดับของข้อมูล ประการที่สองคือการรักษาสัดส่วนของตัวแปรตาม (Stratified Splitting) โดยเฉพาะในปัญหาการจำแนกประเภท หากข้อมูลมีการกระจายไม่สมดุล (Imbalanced) การสุ่มเื่าจะอาจทำให้ชุดฝึกสอนและชุดทดสอบมีสัดส่วนของแต่ละกลุ่มไม่สอดคล้องกัน การใช้ Stratified Splitting จะรักษาสัดส่วนของแต่ละกลุ่มในทั้งสองชุดให้เท่าเดิมกับข้อมูลทั้งหมด

นอกจากการแบ่งข้อมูลเป็นชุดฝึกสอนและชุดทดสอบแล้ว ในทางปฏิบัติมักมีการแบ่งส่วนเพิ่มเติมเรียกว่าชุดตรวจสอบ (Validation Set) ซึ่งใช้สำหรับปรับแต่งพารามิเตอร์และเลือกโมเดล ชุดตรวจสอบจึงมีบทบาทเป็นตัวแทนของข้อมูลใหม่ในขั้นตอนการพัฒนา ทำให้สามารถประเมินและเปรียบเทียบตัวแบบหลายตัวก่อนนำไปทดสอบกับชุดข้อมูลทดสอบจริง การแบ่งข้อมูลทั้งหมดออกเป็นสามส่วน ได้แก่ ชุดฝึกสอน ชุดตรวจสอบ และชุดทดสอบ ในสัดส่วน 60:20:20 หรือ 70:15:15 เป็นแนวทางที่นิยมใช้

\begin{myBox}[]{ข้อควรระวังในการแบ่งข้อมูลฝึกสอนและข้อมูลทดสอบ}{}
การแบ่งข้อมูลอย่างไม่เหมาะสมอาจนำไปสู่การประเมินประสิทธิภาพที่คลาดเคลื่อนอย่างมาก ประเด็นสำคัญที่ต้องระวังมีดังนี้

\textbf{การรั่วไหลของข้อมูล (Data Leakage)} เกิดขึ้นเมื่อข้อมูลจากชุดทดสอบ ``รั่ว'' เข้าสู่กระบวนการฝึกสอน เช่น เมื่อใช้ข้อมูลที่มาจากการคำนวณจากทั้งชุดข้อมูลในการสร้างคุณลักษณะ (Feature) ก่อนการแบ่ง หรือเมื่อมีตัวอย่างเดียวกันปรากฏในทั้งชุดฝึกสอนและชุดทดสอบ (เช่น ข้อมูลลูกค้าคนเดียวกันในช่วงเวลาต่างกัน) การรั่วไหลทำให้ประสิทธิภาพที่วัดได้สูงกว่าความเป็นจริงอย่างมาก

\textbf{การประเมินที่ลำเอียง (Biased Evaluation)} หากชุดทดสอบไม่เป็นตัวแทนที่ดีของข้อมูลทั้งหมด เช่น ในปัญหาการจำแนกประเภทที่มีกลุ่มน้อยมาก (Minority Class) หากกลุ่มนี้ไม่ปรากฏในชุดทดสอบเลย การประเมินจะไม่สะท้อนความสามารถของโมเดลในการจำแนกกลุ่มสำคัญนั้น

\textbf{การ Overfitting กับชุดตรวจสอบ} เมื่อใช้ชุดตรวจสอบในการปรับแต่งพารามิเตอร์หลายครั้งซ้ำๆ อาจเกิดการ Overfitting กับชุดตรวจสอบโดยไม่รู้ตัว ในกรณีนี้ควรใช้เทคนิค Cross-Validation เพื่อประเมินประสิทธิภาพอย่างเป็นกลาง
\end{myBox}

เมื่อพิจารณาชุดข้อมูลที่มีโครงสร้างเป็นกลุ่มหรือมีความสัมพันธ์ภายในกลุ่ม เช่น ข้อมูลผู้ป่วยจากโรงพยาบาลหลายแห่ง หรือข้อมูลจากลูกค้าในหลายสาขาธนาคาร การสุ่มแบ่งข้อมูลโดยไม่คำนึงถึงโครงสร้างกลุ่มอาจทำให้ตัวอย่างจากกลุ่มเดียวกันปรากฏในทั้งชุดฝึกสอนและชุดทดสอบ ทำให้การประเมินประสิทธิภาพไม่แท้จริง ในกรณีเช่นนี้ควรใช้การแบ่งข้อมูลแบบ Group-based Splitting โดยแบ่งตามกลุ่ม เช่น สุ่มแบ่งโรงพยาบาลหรือสาขาธนาคาร ไม่ใช่สุ่มแบ่งผู้ป่วยหรือลูกค้า

สำหรับข้อมูลอนุกรมเวลา (Time Series Data) การสุ่มแบ่งแบบปกติไม่เหมาะสม เนื่องจากข้อมูลมีลำดับเวลาและความสัมพันธ์ระหว่างช่วงเวลาที่ติดกัน ในกรณีนี้ควรใช้การแบ่งแบบ Time-based Splitting โดยใช้ข้อมูลในอดีตสำหรับฝึกสอนและข้อมูลในอนาคตสำหรับทดสอบ ตัวอย่างเช่น หากมีข้อมูลราคาหุ้นย้อนหลัง 5 ปี อาจใช้ข้อมูล 4 ปีแรกสำหรับฝึกสอนและปีสุดท้ายสำหรับทดสอบ การใช้ข้อมูลอนาคตสำหรับฝึกสอนและข้อมูลอดีตสำหรับทดสอบจะทำให้ได้ผลการประเมินที่ไม่สมจริงและไม่มีความหมายในทางปฏิบัติ

\section{ฟังก์ชันการสูญเสีย}
ฟังก์ชันการสูญเสีย (Loss Function) เป็นองค์ประกอบสำคัญในการเรียนรู้แบบมีผู้สอนที่ทำหน้าที่วัดความผิดพลาดหรือความไม่สอดคล้องระหว่างค่าทำนายของตัวแบบกับค่าจริง ฟังก์ชันนี้เป็นตัวชี้วัดที่ใช้ในการปรับพารามิเตอร์ของตัวแบบให้เหมาะสมที่สุด โดยเป้าหมายคือการหาค่าพารามิเตอร์ที่ทำให้ค่าฟังก์ชันการสูญเสียมีค่าน้อยที่สุดเท่าที่จะเป็นไปได้ การเลือกฟังก์ชันการสูญเสียที่เหมาะสมขึ้นอยู่กับลักษณะของปัญหา ทั้งประเภทของตัวแปรตามและคุณสมบัติที่ต้องการของตัวประมาณ

สำหรับปัญหาการพยากรณ์หรือการถดถอย ฟังก์ชันการสูญเสียที่นิยมใช้คือ Mean Squared Error (MSE) ซึ่งมีสูตรดังนี้

\[
\mathcal{L}_{\text{MSE}}(y, \hat{y}) = \frac{1}{n}\sum_{i=1}^{n}(y_i - \hat{y}_i)^{2}
\]

โดยที่ $y_i$ คือค่าจริงของตัวอย่างที่ $i$ และ $\hat{y}_i$ คือค่าทำนายของตัวแบบสำหรับตัวอย่างเดียวกัน การยกกำลังสองทำให้ค่าความผิดพลาดที่มากขึ้นมีผลกระทบมากกว่าค่าความผิดพลาดที่น้อย ซึ่งเหมาะสมในสถานการณ์ที่ต้องการลดค่าความผิดพลาดใหญ่ให้น้อยที่สุด อย่างไรก็ตาม ข้อเสียของ MSE คือความไวต่อค่าผิดปกติ (Outliers) เนื่องจากการยกกำลังสองจะขยายผลของค่าผิดปกติมาก

ฟังก์ชันการสูญเสียอีกประเภทหนึ่งสำหรับการพยากรณ์คือ Mean Absolute Error (MAE) ซึ่งมีสูตรดังนี้

\[
\mathcal{L}_{\text{MAE}}(y, \hat{y}) = \frac{1}{n}\sum_{i=1}^{n}|y_i - \hat{y}_i|
\]

MAE มีความไวต่อค่าผิดปกติน้อยกว่า MSE เนื่องจากใช้ค่าสัมบูรณ์แทนการยกกำลังสอง การเลือกระหว่าง MSE และ MAE ขึ้นอยู่กับลักษณะของข้อมูลและว่าต้องการให้ตัวแบบมีความไวต่อค่าผิดปกติมากน้อยเพียงใด นอกจากนี้ยังมี Huber Loss ซึ่งเป็นการผสมผสานระหว่าง MSE และ MAE โดยใช้ MSE สำหรับค่าความผิดพลาดเล็กๆ และใช้ MAE สำหรับค่าความผิดพลาดใหญ่ ทำให้ได้ประโยชน์จากทั้งสองแบบ

สำหรับปัญหาการจำแนกประเภทแบบทวิภาค ฟังก์ชันการสูญเสียที่นิยมใช้คือ Cross-Entropy Loss หรือที่เรียกว่า Binary Cross-Entropy ซึ่งมีสูตรดังนี้

\[
\mathcal{L}_{\text{BCE}}(y, \hat{p}) = -\frac{1}{n}\sum_{i=1}^{n}\left[y_i \log(\hat{p}_i) + (1 - y_i)\log(1 - \hat{p}_i)\right]
\]

โดยที่ $y_i \in \{0, 1\}$ คือค่าฉลากจริงและ $\hat{p}_i = P(y_i = 1|\mathbf{x}_i)$ คือความน่าจะเป็นที่ตัวแบบทำนายว่าตัวอย่างที่ $i$ เป็นกลุ่ม 1 ฟังก์ชันนี้มีรากฐานจากทฤษฎีสารสนเทศ (Information Theory) และมีคุณสมบัติที่เหมาะสมสำหรับการประมาณความน่าจะเป็น การใช้ Cross-Entropy แทน 0-1 Loss (ซึ่งให้ค่า 0 หากทำนายถูกและ 1 หากทำนายผิด) ช่วยให้การหาค่าเหมาะสมที่สุดมีความลื่นไหล (Smooth) และสามารถใช้ Gradient-based Optimization ได้อย่างมีประสิทธิภาพ

ในกรณีของการจำแนกประเภทหลายกลุ่ม ฟังก์ชัน Cross-Entropy ขยายไปสู่ Multi-class Cross-Entropy ซึ่งมีสูตรดังนี้

\[
\mathcal{L}_{\text{CE}}(y, \hat{\mathbf{p}}) = -\sum_{i=1}^{n}\sum_{k=1}^{K}y_{ik}\log(\hat{p}_{ik})
\]

โดยที่ $K$ คือจำนวนกลุ่ม $y_{ik} = 1$ หากตัวอย่างที่ $i$ เป็นสมาชิกของกลุ่ม $k$ และ $\hat{p}_{ik}$ คือความน่าจะเป็นที่ตัวแบบทำนายว่าตัวอย่างที่ $i$ เป็นสมาชิกของกลุ่ม $k$ ฟังก์ชันนี้เป็นพื้นฐานของการเรียนรู้ของโครงข่ายประสาทเทียม (Neural Networks) สำหรับงานจำแนกประเภทหลายกลุ่ม

\begin{myBox}[]{ตัวอย่างการคำนวณฟังก์ชันการสูญเสีย}{}
เพื่อให้เข้าใจการทำงานของฟังก์ชันการสูญเสีย พิจารณาตัวอย่างการพยากรณ์ราคาบ้านด้วย MSE

สมมติมีข้อมูลราคาบ้านจริง 3 หลัง ได้แก่ $y = [3.5, 5.2, 2.8]$ ล้านบาท และค่าทำนายจากตัวแบบ $\hat{y} = [3.2, 5.5, 2.9]$ ล้านบาท ค่า MSE คำนวณได้ดังนี้

\[
\text{MSE} = \frac{1}{3}\left[(3.5-3.2)^2 + (5.2-5.5)^2 + (2.8-2.9)^2\right] = \frac{1}{3}\left[0.09 + 0.09 + 0.01\right] = \frac{0.19}{3} \approx 0.063
\]

สำหรับการจำแนกประเภท พิจารณาข้อมูลอีเมล 3 ฉบับ โดยฉลากจริง $y = [1, 0, 1]$ (1=สแปม 0=ไม่ใช่สแปม) และความน่าจะเป็นที่ตัวแบบทำนาย $\hat{p} = [0.9, 0.3, 0.7]$ ค่า BCE คำนวณได้ดังนี้

\[
\text{BCE} = -\frac{1}{3}\left[\log(0.9) + \log(0.7) + \log(1-0.3)\right] = -\frac{1}{3}\left[-0.105 + (-0.357) + (-0.357)\right] = 0.273
\]

ค่า BCE ยิ่งต่ำ หมายถึงตัวแบบยิ่งทำนายได้แม่นยำ
\end{myBox}

\begin{exmp}
\label{exmp:mae_vs_mse_comparison}
เพื่อเปรียบเทียบคุณสมบัติของ MSE และ MAE พิจารณาข้อมูลการพยากรณ์ยอดขายรายวัน 5 วัน ดังตารางต่อไปนี้

\begin{table}[h!]
\centering
\caption{ข้อมูลเปรียบเทียบ MSE กับ MAE}
\label{tab:mse_mae_comparison}
\begin{tabular}{|c|c|c|c|c|}
\hline
\textbf{วัน} & \textbf{ยอดขายจริง (พันบาท)} & \textbf{ค่าทำนาย (พันบาท)} & \textbf{$|y_i - \hat{y}_i|$} & \textbf{$(y_i - \hat{y}_i)^2$} \\ \hline
1 & 10.0 & 10.5 & 0.5 & 0.25 \\ \hline
2 & 12.0 & 12.2 & 0.2 & 0.04 \\ \hline
3 & 11.5 & 11.0 & 0.5 & 0.25 \\ \hline
4 & 13.0 & 13.8 & 0.8 & 0.64 \\ \hline
5 & 10.5 & 10.4 & 0.1 & 0.01 \\ \hline
\end{tabular}
\end{table}

จากข้อมูลในตาราง ค่า MAE และ MSE คำนวณได้ดังนี้

\[
\text{MAE} = \frac{0.5 + 0.2 + 0.5 + 0.8 + 0.1}{5} = \frac{2.1}{5} = 0.42
\]

\[
\text{MSE} = \frac{0.25 + 0.04 + 0.25 + 0.64 + 0.01}{5} = \frac{1.19}{5} = 0.238
\]

สังเกตได้ว่าวันที่ 4 ซึ่งมีความผิดพลาดมากที่สุด (0.8 พันบาท) มีสัดส่วนต่อ MSE สูงมากเมื่อเทียบกับ MAE เนื่องจากการยกกำลังสองจะขยายผลของความผิดพลาดที่ใหญ่ ดังนั้นในสถานการณ์ที่มีค่าผิดปกติ (Outliers) ควรพิจารณาใช้ MAE แทน MSE เพื่อให้ได้ตัวแบบที่ทนทานต่อค่าผิดปกติมากขึ้น
\end{exmp}

นอกจากฟังก์ชันการสูญเสียสำหรับการพยากรณ์และการจำแนกประเภทแบบมาตรฐานแล้ว ยังมีฟังก์ชันการสูญเสียเฉพาะทางสำหรับปัญหาเฉพาะประเภท เช่น Focal Loss ซึ่งออกแบบมาสำหรับปัญหาที่มีการกระจายไม่สมดุลรุนแรง (Class Imbalance) โดยลดผลกระทบของตัวอย่างที่ทำนายง่าย (Easy Examples) และให้ความสำคัญกับตัวอย่างที่ทำนายยาก (Hard Examples) หรือ Huber Loss ซึ่งเป็นการผสมระหว่าง MSE และ MAE ตามที่กล่าวไว้ข้างต้น เพื่อให้ได้ความทนทานต่อค่าผิดปกติในขณะที่ยังคงความลื่นไหลในการหาค่าเหมาะสมที่สุด

การเลือกฟังก์ชันการสูญเสียที่เหมาะสมมีผลอย่างมากต่อประสิทธิภาพของตัวแบบ ฟังก์ชันการสูญเสียที่ดีควรสะท้อนต้นทุนที่แท้จริงของความผิดพลาดในบริบทของปัญหา ตัวอย่างเช่น ในระบบวินิจฉัยโรค การทำนายผลบวกลวง (False Positive) อาจมีต้นทุนต่างจากการทำนายผลลบลวง (False Negative) ดังนั้นอาจต้องใช้ฟังก์ชันการสูญเสียที่ให้น้ำหนักต่างกันกับความผิดพลาดแต่ละประเภท หรืออาจปรับเปลี่ยนฟังก์ชันการสูญเสียมาตรฐานด้วยการกำหนดค่าน้ำหนัก (Class Weights) เพื่อสะท้อนความสำคัญที่แตกต่างกันของการทำนายผิดแต่ละประเภท

\section{การหาค่าเหมาะสมที่สุด}
การหาค่าเหมาะสมที่สุด (Optimization) เป็นกระบวนการปรับพารามิเตอร์ของตัวแบบให้ลดค่าฟังก์ชันการสูญเสียให้น้อยที่สุดบนข้อมูลฝึกสอน หรือเพิ่มค่าฟังก์ชันที่วัดความแม่นยำให้มากที่สุดในกรณีของบางตัวแบบ ในทางคณิตศาสตร์ ปัญหาการหาค่าเหมาะสมที่สุดสามารถเขียนได้ว่าต้องการหาค่าพารามิเตอร์ $\boldsymbol{\theta}^{*}$ ที่

\[
\boldsymbol{\theta}^{*} = \underset{\boldsymbol{\theta}}{\arg\min} \; \mathcal{L}(\boldsymbol{\theta}) = \underset{\boldsymbol{\theta}}{\arg\min} \; \frac{1}{n}\sum_{i=1}^{n}\ell\!\left(y_i,\; f(\mathbf{x}_i; \boldsymbol{\theta})\right)
\]

โดยที่ $\mathcal{L}(\boldsymbol{\theta})$ คือฟังก์ชันการสูญเสียรวมบนข้อมูลฝึกสอนทั้งหมด $\ell(\cdot)$ คือฟังก์ชันสูญเสียรายตัวอย่าง $f(\mathbf{x}; \boldsymbol{\theta})$ คือตัวแบบ และ $n$ คือจำนวนตัวอย่างในข้อมูลฝึกสอน การหาค่าเหมาะสมที่สุดของ $\boldsymbol{\theta}$ ให้ได้ค่าฟังก์ชันการสูญเสียต่ำที่สุดทำได้หลายวิธี ซึ่งแต่ละวิธีมีข้อดีข้อเสียและความเหมาะสมต่างกันสำหรับปัญหาแต่ละประเภท

วิธีการหาค่าเหมาะสมที่สุดแบบเชิงวิเคราะห์ (Analytical Solution) ใช้ได้เฉพาะกับตัวแบบที่มีรูปแบบง่ายและฟังก์ชันการสูญเสียที่มีผลเฉลยตรง (Closed-form Solution) ตัวอย่างที่เห็นได้ชัดคือการถดถอยเชิงเส้นอย่างง่าย (Simple Linear Regression) ที่ใช้วิธีกำลังสองน้อยที่สุด (Ordinary Least Squares) ซึ่งมีผลเฉลยตรงดังนี้

\[
\boldsymbol{\hat{\beta}} = (\mathbf{X}^{T}\mathbf{X})^{-1}\mathbf{X}^{T}\mathbf{y}
\]

อย่างไรก็ตาม ตัวแบบส่วนใหญ่ในการเรียนรู้ของเครื่องไม่มีผลเฉลยตรง หรือการคำนวณผลเฉลยตรงมีความซับซ้อนเกินไปสำหรับข้อมูลขนาดใหญ่ จึงต้องใช้วิธีการหาค่าเหมาะสมที่สุดเชิงตัวเลข (Numerical Optimization) แทน

วิธีการหาค่าเหมาะสมที่สุดเชิงตัวเลขที่พื้นฐานที่สุดคือ Gradient Descent ซึ่งอาศัยอนุพันธ์ของฟังก์ชันการสูญเสียเพื่อปรับพารามิเตอร์ไปในทิศทางที่ลดค่าฟังก์ชันการสูญเสียมากที่สุด หลักการคือ หากฟังก์ชันการสูญเสีย $\mathcal{L}(\boldsymbol{\theta})$ มีค่าลดลงเมื่อเคลื่อนที่ไปในทิศทางที่ตรงข้ามกับ Gradient $\nabla\mathcal{L}(\boldsymbol{\theta})$ ดังนั้นการอัปเดตพารามิเตอร์ตามสูตร

\[
\boldsymbol{\theta} \leftarrow \boldsymbol{\theta} - \alpha \nabla\mathcal{L}(\boldsymbol{\theta})
\]

จะทำให้ค่าฟังก์ชันการสูญเสียลดลง โดยที่ $\alpha > 0$ คืออัตราการเรียนรู้ (Learning Rate) ซึ่งควบคุมขนาดของก้าวในแต่ละการอัปเดต การเลือกอัตราการเรียนรู้ที่เหมาะสมมีความสำคัญมาก หากมากเกินไป การอัปเดตอาจกระโดดข้ามค่าต่ำสุดและทำให้ไม่ converges หากน้อยเกินไป การ convergence จะช้ามากและอาจติดอยู่ที่ค่าต่ำสุดเฉพาะที่ (Local Minimum)

\begin{myBox}[]{วิธีการหาค่าเหมาะสมที่สุดแบบต่างๆ}{}
วิธีการหาค่าเหมาะสมที่สุดสำหรับการเรียนรู้ของเครื่องสามารถจำแนกได้หลายประเภท ดังนี้

\textbf{Gradient Descent แบบเต็ม (Batch Gradient Descent)} คำนวณ Gradient จากข้อมูลฝึกสอนทั้งหมดในแต่ละการอัปเดต วิธีนี้ให้การประมาณค่า Gradient ที่แม่นยำแต่ช้าสำหรับข้อมูลขนาดใหญ่

\textbf{Stochastic Gradient Descent (SGD)} คำนวณ Gradient จากตัวอย่างเดียวหรือกลุ่มเล็กๆ ในแต่ละการอัปเดต วิธีนี้เร็วกว่าและสามารถหลบหนีจากค่าต่ำสุดเฉพาะที่ได้ดีกว่า แต่การประมาณค่า Gradient มีความแปรปรวนสูง

\textbf{Mini-batch Gradient Descent} เป็นการผสมผสานระหว่าง Batch และ Stochastic โดยใช้กลุ่มย่อย (Mini-batch) ของตัวอย่างในแต่ละการอัปเดต วิธีนี้ได้ประโยชน์จากทั้งสองแบบ คือความเร็วและความเสถียร

\textbf{Adaptive Optimization Methods} เช่น Adam, RMSprop, AdaGrad ปรับอัตราการเรียนรู้สำหรับแต่ละพารามิเตอร์โดยอัตโนมัติตามประวัติของ Gradient ทำให้การหาค่าเหมาะสมที่สุดมีประสิทธิภาพมากขึ้นโดยเฉพาะสำหรับโครงข่ายประสาทเทียม
\end{myBox}

\begin{exmp}
\label{exmp:gradient_descent_simple}
เพื่อให้เข้าใจการทำงานของ Gradient Descent อย่างเป็นรูปธรรม พิจารณาฟังก์ชันการสูญเสียแบบง่ายที่มีตัวแปร $\theta$ เพียงตัวเดียว

\[
\mathcal{L}(\theta) = (\theta - 3)^2
\]

เป้าหมายคือการหาค่า $\theta$ ที่ทำให้ $\mathcal{L}(\theta)$ มีค่าน้อยที่สุด จากสมการจะเห็นได้ว่าค่าต่ำสุดอยู่ที่ $\theta^* = 3$ เมื่อ $\mathcal{L}(3) = 0$

สมมติเริ่มต้นที่ $\theta_0 = 0$ และกำหนดอัตราการเรียนรา $\alpha = 0.1$ Gradient ของ $\mathcal{L}(\theta)$ คือ

\[
\frac{d\mathcal{L}}{d\theta} = 2(\theta - 3)
\]

\textbf{รอบที่ 1:} ที่ $\theta_0 = 0$ จะได้ $\nabla\mathcal{L} = 2(0-3) = -6$

\[
\theta_1 = \theta_0 - \alpha \nabla\mathcal{L} = 0 - 0.1(-6) = 0.6
\]

\textbf{รอบที่ 2:} ที่ $\theta_1 = 0.6$ จะได้ $\nabla\mathcal{L} = 2(0.6-3) = -4.8$

\[
\theta_2 = 0.6 - 0.1(-4.8) = 1.08
\]

\textbf{รอบที่ 3:} ที่ $\theta_2 = 1.08$ จะได้ $\nabla\mathcal{L} = 2(1.08-3) = -3.84$

\[
\theta_3 = 1.08 - 0.1(-3.84) = 1.464
\]

สังเกตได้ว่าค่า $\theta$ ค่อย ๆ ก้าวเข้าใกล้ 3 จาก 0 $\to$ 0.6 $\to$ 1.08 $\to$ 1.464 ซึ่งเป็นไปในทิศทางที่ลดค่าฟังก์ชันการสูญเสีย หากทำต่อไปอีกหลายรอบ ค่า $\theta$ จะ converge ไปยัง $\theta^* = 3$
\end{exmp}

ในการใช้งานจริง มักใช้ Mini-batch Gradient Descent ร่วมกับ Adaptive Optimization Methods เช่น Adam ซึ่งเป็นวิธีที่ได้รับความนิยมมากที่สุดในปัจจุบันสำหรับการฝึกโครงข่ายประสาทเทียม Adam ปรับอัตราการเรียนรู้สำหรับแต่ละพารามิเตอร์โดยอาศัยทั้ง First Moment (Mean ของ Gradient) และ Second Moment (Variance ของ Gradient) ทำให้การ convergence เร็วและมีเสถียรภาพมากกว่า SGD แบบธรรมดา

นอกจากการหาค่าเหมาะสมที่สุดบนข้อมูลฝึกสอนแล้ว ในทางปฏิบัติต้องคำนึงถึงปรากฏการณ์ Overfitting ซึ่งเกิดขึ้นเมื่อตัวแบบมีความซับซ้อนมากเกินไปจนสามารถจดจำข้อมูลฝึกสอนได้อย่างสมบูรณ์แบบแต่ไม่สามารถทำนายข้อมูลใหม่ได้ดี เพื่อแก้ปัญหานี้ มักเพิ่มเทอม Regularization เข้าไปในฟังก์ชันการสูญเสีย เช่น L1 Regularization (Lasso) หรือ L2 Regularization (Ridge) ซึ่งทำหน้าที่ลงโทษค่าพารามิเตอร์ที่มากเกินไป ทำให้ตัวแบบมีความซับซ้อนลดลงและมีความสามารถในการ generalize ไปยังข้อมูลใหม่ดีขึ้น

การหาค่าเหมาะสมที่สุดในปัญหาการเรียนรู้ของเครื่องมีความท้าทายหลายประการ ประการแรกคือปัญหา Local Minimum เนื่องจากฟังก์ชันการสูญเสียมักมีความไม่เป็นเชิงเส้นสูง อาจมีค่าต่ำสุดเฉพาะที่หลายค่า อัลกอริทึม Gradient Descent อาจติดอยู่ที่ค่าต่ำสุดเฉพาะที่แทนที่จะเป็นค่าต่ำสุดสมบูรณ์ (Global Minimum) อย่างไรก็ตาม งานวิจัยล่าสุดพบว่าสำหรับโครงข่ายประสาทเทียมขนาดใหญ่ ค่าต่ำสุดเฉพาะที่มักมีค่าใกล้เคียงกันมากพอที่จะไม่เป็นปัญหาในทางปฏิบัติ ประการที่สองคือปัญหา Vanishing และ Exploding Gradient ซึ่งเกิดขึ้นเมื่อ Gradient มีค่าน้อยมากจนไม่สามารถอัปเดตพารามิเตอร์ได้ หรือมีค่ามากจนทำให้การอัปเดตผิดปกติรุนแรง ปัญหาเหล่านี้ได้รับการแก้ไขบางส่วนด้วยเทคนิคต่างๆ เช่น Batch Normalization และ Residual Connections

\newpage
%\RefChapter
